% Cote 119: offprints, reports, courses, CV and letters (typed and printed part): archive cote 119, site page 234 (PDF page 235).
% typescript. Draft compared: none.
% AI-transcribed by Claude Opus 5.5 (claude-opus-5-5, adaptive thinking, effort high), Message Batches API, batch r119_full (msgbatch_01DwZNm5Y7YsKBWyAN4sFs8w), 2026-11-13; not yet checked by eye.
\scanpage{119/234}
\mspage{- 21 -}

Complément à la bibliographie sommaire jointe au
Curriculum Vitae de A. Grothendieck (travaux non inclus
dans ladite bibliographie).

\uline{Analyse fonctionnelle}.

[1 bis] Sur les espaces de solutions d'une classe
générale d'équations aux dérivées partielles,
Journal d'Analyse Math. vol. II, pp. 243-280
(1952/53).

[2 bis] Sur certaines classes de suites dans les espaces
de Banach, et le théorème de Dvoretzky-Rogers,
Boletim da Soc. Mat. de Sa\unsure{\~o} Paulo, Vol. 8°,
pp. 85-110 (1953).

[3 bis] Sur les applications linéaires faiblement
compactes d'espaces du type $C(K)$, Canadian
Journal of Math., Vol. 5, pp. 129-173 (1953).

[4 bis] Sur certains sous-espaces vectoriels de $L^p$ ,
Can. J. Math. vol. 6, pp. 158-160 (1953).

[5 bis] Une caractérisation vectorielle-métrique des
espaces $L^1$ , Can. Journ. Math. vol. 7,
pp. 552-561 (1955).

[6 bis] Réarrangements de fonctions et inégalités de
convexité dans les algèbres de von Neumann
munies d'une trace, Séminaire Bourbaki n° 115
(Mars 1955).

[7 bis] Un résultat sur le dual d'une $C^{\unsure{\hand{$*$}}}$-algèbre,
Journ. de Math. vol. 36, pp. 97-108 (1957).

[8 bis] The trace of certain operators, Studia
Mathematica t. 20 (1961) pp. 141-143.

\margin{234}
